%----------------------------------------------------------------------------------------
% VECTOREN IN DE RUIMTE. Alle opgaven zijn eigen werk.
%
% Opbouw: Start, leerstof (theorie, werkwijze of uitleg, voorbeeld, korte oefening), en op een
% nieuwe pagina de Opdrachten (door elkaar, stijgende moeilijkheid), Koppel en Test jezelf.
% Labels van dit hoofdstuk beginnen met h1-.
%----------------------------------------------------------------------------------------

\mychapter{1}{Vectoren in de ruimte}

\lkr{Leerplan B+S'': LPD 5-8 (ruimtemeetkunde), het deel over coördinaten en vectoren in de ruimte. Leerlingen kennen vectoren en coördinaten in het vlak uit de 2de graad; hier komt er een derde coördinaat bij. Laat ze de figuren zelf tekenen in schuine projectie, en controleer met een kubus of balk: dat traint het ruimtelijk inzicht.}

Een drone hangt $3$ meter boven een punt van een sportveld. Om zijn plaats te geven, heb je drie getallen nodig: twee voor het punt op het veld, en de hoogte. Met drie coördinaten en met vectoren reken je in de ruimte zoals je in het vlak gewoon bent.

\begin{oefeningen}[Start]
\oef{h1-start-1}

\oef{h1-verken}
\end{oefeningen}

%========================================================================================
\section{Coördinaten in de ruimte}

\naast{%
Een \textbf{orthonormaal assenstelsel} in de ruimte heeft drie onderling loodrechte assen door de oorsprong $O$, met dezelfde lengte-eenheid. Een punt $P$ heeft drie coördinaten $(x, y, z)$: je bereikt het vanuit $O$ door $x$ langs de $x$-as, $y$ langs de $y$-as en $z$ omhoog te gaan.

In de figuur ligt de balk met hoekpunten $O$ en $P(2, 3, 2)$ met zijn ribben langs de assen.}{%
\begin{tikzpicture}[ruimte,scale=.85,font=\scriptsize]
  \assendrie{4}
  \draw[accent] (2,0,0) -- (2,3,0) -- (0,3,0) (2,0,0) -- (2,0,2) -- (0,0,2) -- (0,3,2) -- (2,3,2) -- (2,0,2) (2,3,0) -- (2,3,2) (0,3,0) -- (0,3,2);
  \fill[accent] (2,3,2) circle (.06) node[right,text=black] {$P(2, 3, 2)$};
  \node[left] at (2,0,0) {$2$}; \node[below] at (0,3,0) {$3$}; \node[left] at (0,0,2) {$2$};
\end{tikzpicture}}

\begin{oefeningen}
\oef{h1-coordinaten}
\end{oefeningen}

\begin{theorie}
De \textbf{\textcolor{accent}{afstand}} tussen $A(x_1, y_1, z_1)$ en $B(x_2, y_2, z_2)$ en het \textbf{\textcolor{accent}{midden}} van $[AB]$ zijn
\[\formule{|AB| = \sqrt{(x_2 - x_1)^2 + (y_2 - y_1)^2 + (z_2 - z_1)^2}} \qquad \formule{M\Bigl(\tfrac{x_1 + x_2}{2}, \tfrac{y_1 + y_2}{2}, \tfrac{z_1 + z_2}{2}\Bigr)}.\]
\end{theorie}

De afstandsformule is twee keer de stelling van Pythagoras: eerst in het horizontale vlak, dan in de verticale driehoek met de hoogte.

\voorbeeld{}{Bereken de afstand tussen $A(1, 2, 3)$ en $B(4, -2, 15)$, en het midden van $[AB]$.\label{vb:h1-afstand}}{%
\uitleg{Neem de verschillen van de coördinaten.}
\begin{stappenbreed}
|AB| = \sqrt{3^2 + (-4)^2 + 12^2} = \sqrt{169} = 13 & Verschillen $3$, $-4$, $12$. \\
M\bigl(\tfrac52, 0, 9\bigr) & Gemiddelden.
\end{stappenbreed}}

\begin{oefeningen}
\oef{h1-afstand}

\oef{h1-midden}
\end{oefeningen}

%========================================================================================
\needspace{20\baselineskip}
\section{Vectoren}

\begin{theorie}
De \textbf{\textcolor{accent}{vector}} $\overrightarrow{AB}$ van $A(x_1, y_1, z_1)$ naar $B(x_2, y_2, z_2)$ heeft als coördinaten de verschillen:
\[\formule{\overrightarrow{AB} = \vc{x_2 - x_1}{y_2 - y_1}{z_2 - z_1}}.\]
Vectoren tel je op en vermenigvuldig je met een getal per coördinaat. De \textbf{\textcolor{accent}{norm}} (lengte) van $\vec v = \vcr{a}{b}{c}$ is $\formule{\norm{\vec v} = \sqrt{a^2 + b^2 + c^2}}$.
\end{theorie}

Een vector is een verplaatsing: hij heeft een richting, een zin en een lengte, maar geen vaste plaats. Twee pijlen met dezelfde richting, zin en lengte stellen dezelfde vector voor. Zo is $B = A + \overrightarrow{AB}$: vertrek in $A$ en verplaats.

\voorbeeld{}{Gegeven $A(1, 0, 2)$, $B(3, 4, -2)$ en $C(0, 1, 1)$. Bereken $\overrightarrow{AB}$, $\norm{\overrightarrow{AB}}$ en het punt $D$ waarvoor $ABDC$ een parallellogram is.\label{vb:h1-vectoren}}{%
\uitleg{In een parallellogram $ABDC$ is $\overrightarrow{CD} = \overrightarrow{AB}$: $D = C + \overrightarrow{AB}$.}
\begin{stappenbreed}
\overrightarrow{AB} = \vc{2}{4}{-4},\quad \norm{\overrightarrow{AB}} = \sqrt{4 + 16 + 16} = 6 & \\
D = C + \overrightarrow{AB} = (0 + 2, 1 + 4, 1 - 4) = (2, 5, -3) &
\end{stappenbreed}}

\begin{oefeningen}
\oef{h1-vectoren}

\oef{h1-parallellogram}

\oef{h1-norm}
\end{oefeningen}

%========================================================================================
\needspace{12\baselineskip}
\section{Collineaire en coplanaire vectoren}

\begin{theorie}
Twee vectoren $\vec v$ en $\vec w \neq \vec 0$ zijn \textbf{\textcolor{accent}{collineair}} (evenwijdig) als $\formule{\vec v = k\,\vec w}$ voor een getal $k$.

Drie vectoren zijn \textbf{\textcolor{accent}{coplanair}} als ze in één vlak liggen: als een van hen een lineaire combinatie van de andere twee is, $\formule{\vec u = r\,\vec v + s\,\vec w}$.
\end{theorie}

Drie punten $A$, $B$, $C$ liggen op één rechte als $\overrightarrow{AB}$ en $\overrightarrow{AC}$ collineair zijn. Vier punten liggen in één vlak als $\overrightarrow{AB}$, $\overrightarrow{AC}$ en $\overrightarrow{AD}$ coplanair zijn.

\voorbeeld{}{Liggen $A(1, 2, 3)$, $B(2, 4, 7)$ en $C(4, 8, 15)$ op één rechte?\label{vb:h1-collineair}}{%
\uitleg{Vergelijk $\overrightarrow{AB}$ en $\overrightarrow{AC}$: is de ene een veelvoud van de andere?}
\begin{stappenbreed}
\overrightarrow{AB} = \vc{1}{2}{4},\quad \overrightarrow{AC} = \vc{3}{6}{12} = 3\,\overrightarrow{AB} & Ja: collineair, dus op één rechte.
\end{stappenbreed}}

\begin{oefeningen}
\oef{h1-collineair}
\end{oefeningen}

\voorbeeld{}{Liggen $O(0, 0, 0)$, $B(1, 0, 1)$, $C(0, 1, 1)$ en $D(2, 3, 4)$ in één vlak?\label{vb:h1-coplanair}}{%
\uitleg{Zoek $r$ en $s$ met $\overrightarrow{OD} = r\,\overrightarrow{OB} + s\,\overrightarrow{OC}$: een stelsel van drie vergelijkingen met twee onbekenden.}
\begin{stappenbreed}
\vc{2}{3}{4} = r\vc{1}{0}{1} + s\vc{0}{1}{1} & \\
r = 2,\ s = 3,\ \text{controle } r + s = 5 \neq 4 & Derde vergelijking klopt niet.
\end{stappenbreed}
Er is geen oplossing: de vier punten liggen \textit{niet} in één vlak. Met $D(2, 3, 5)$ zou het wel zo zijn.}

\begin{werkwijze}
Collineair: kijk of de verhoudingen van de coördinaten gelijk zijn (let op nullen). Coplanair: schrijf $\vec u = r\vec v + s\vec w$ per coördinaat uit, los twee vergelijkingen op naar $r$ en $s$, en controleer met de derde.
\end{werkwijze}

\begin{oefeningen}
\oef{h1-coplanair}
\end{oefeningen}

%========================================================================================
\opdrachtenkop

\begin{oefeningen}
\oef{h1-o13}

\needspace{12\baselineskip}
\oef{h1-o9}

\needspace{10\baselineskip}
\oef{h1-o14}

\oef{h1-o1}

\oef{h1-o12}

\needspace{14\baselineskip}
\oef{h1-o10}

\oef{h1-o15}

\oef{h1-o3}

\oef{h1-o2}

\needspace{14\baselineskip}
\oef{h1-o11}

\oef{h1-o4}

\oef{h1-o19}

\oef{h1-o16}

\needspace{14\baselineskip}
\oef{h1-o22}

\oef{h1-o18}

\oef{h1-o6}

\needspace{14\baselineskip}
\oef{h1-o17}

\oef{h1-o21}

\oef{h1-o5}

\oef{h1-o23}

\oef{h1-o20}

\oef{h1-o26}

\needspace{16\baselineskip}
\oef{h1-o25}

\oef{h1-o7}

\oef{h1-o24}

\oef{h1-o8}

\oef{h1-o27}
\end{oefeningen}

%----------------------------------------------------------------------------------------
% KOPPEL: verbanden met andere leerstof
%----------------------------------------------------------------------------------------
\needspace{20\baselineskip}
\opdrachtsectie{Koppel}

\lkr{Matrices, Telproblemen en Complexe getallen: kijk na wat al gezien is. Het pakket legt de volgorde niet vast.}

\begin{oefeningen}
\oef{h1-koppel-1}

\oef{h1-koppel-2}

\oef{h1-koppel-3}
\end{oefeningen}

%----------------------------------------------------------------------------------------
\opdrachtsectie{Test jezelf}

\begin{oefeningen}
\oef{h1-t5}

\oef{h1-t1}

\oef{h1-t6}

\oef{h1-t2}

\oef{h1-t7}

\oef{h1-t3}

\oef{h1-t4}
\end{oefeningen}

\herhaalschema{%
  \fzv{oef:h1-t5}{oef:h1-coordinaten}%
  \fv{oef:h1-t1}{vb:h1-afstand}{oef:h1-afstand}%
  \fv{oef:h1-t6}{vb:h1-afstand}{oef:h1-midden}%
  \fv{oef:h1-t2}{vb:h1-vectoren}{oef:h1-vectoren}%
  \fv{oef:h1-t7}{vb:h1-vectoren}{oef:h1-norm}%
  \fv{oef:h1-t3}{vb:h1-collineair}{oef:h1-collineair}%
  \fv{oef:h1-t4}{vb:h1-coplanair}{oef:h1-coplanair}}

\ifzelfstudie
\section*{Maak je eigen samenvatting}
Sluit dit boek en schrijf in je schrift: hoe je een punt in een ruimtelijk assenstelsel tekent, de formules voor afstand en midden, de coördinaten van $\overrightarrow{AB}$, de norm, en wanneer vectoren collineair of coplanair zijn. Vergelijk daarna met de kaders.
\fi
